% GENERATED FILE. Edit canonical lessons, then run npm run generate:books.
% canonical-source-hash: fnv1a64:0d5af1e4d819b673
% canonical-lessons: KA-C165-kadi, KA-C165-nenapisikollu, KA-C165-urulu, KA-C165-uli, KA-C165-kayisu

\chapter{To steal, To recall, To roll, To remain, To heat}
\label{ch:ka-to-steal-to-recall-to-roll-to-remain-to-heat}
\hlchaptermodality{165}

\begin{chapteropening}
\textbf{By the end of this chapter:} \emph{I can say to steal, to recall, to roll, to remain and to heat.}

Complete the last lesson of chapter 165: I can say to steal, to recall, to roll, to remain and to heat.
\end{chapteropening}

\section[\texorpdfstring{\kn{ಕದಿ} (kadi)}{ಕದಿ (kadi)}]{\kn{ಕದಿ} (kadi) — to steal}

\label{lesson:KA-C165-kadi}



\begin{quote}
Before the new one: say the Kannada for to switch on, then the Kannada for to switch off.
\end{quote}

\subsection*{You'll want to know: \kn{ಕದಿ}}
\textbf{\kn{ಕದಿ}} — \emph{kadi} — \textquotedblleft{}to steal\textquotedblright{}.

\begin{grammarlens}[title={What you've built}]
One more word for this chapter.
\end{grammarlens}

\subsection*{Your turn}
\begin{itemize}
\raggedright
  \item \emph{Say it:} \emph{kadi}
  \item \emph{Say it:} \emph{kadi}, once more
  \item \emph{Say it:} say \emph{kadi}
  \item \emph{Recall:} say the Kannada for to switch on, then the Kannada for to switch off, then say \emph{kadi} again
\end{itemize}

\begin{tcolorbox}[breakable,colback=teal!4,colframe=teal!35!black,title={Before you move on}]
What does \kn{ಕದಿ} mean? (\textquotedblleft{}To steal\textquotedblright{}.) Say it once more.
\end{tcolorbox}
\section[\texorpdfstring{\kn{ನೆನಪಿಸಿಕೊಳ್ಳು} (nenapisikoḷḷu)}{ನೆನಪಿಸಿಕೊಳ್ಳು (nenapisikoḷḷu)}]{\kn{ನೆನಪಿಸಿಕೊಳ್ಳು} (nenapisikoḷḷu) — to recall, to remember}

\label{lesson:KA-C165-nenapisikollu}



\begin{quote}
Before the new one: say the Kannada for to switch off, then the Kannada for to steal.
\end{quote}

\subsection*{You'll want to know: \kn{ನೆನಪಿಸಿಕೊಳ್ಳು}}
\textbf{\kn{ನೆನಪಿಸಿಕೊಳ್ಳು}} — \emph{nenapisikoḷḷu} — \textquotedblleft{}to recall, to remember\textquotedblright{}.

\begin{grammarlens}[title={What you've built}]
One more word for this chapter.
\end{grammarlens}

\subsection*{Your turn}
\begin{itemize}
\raggedright
  \item \emph{Say it:} \emph{nenapisikoḷḷu}
  \item \emph{Say it:} \emph{nenapisikoḷḷu}, once more
  \item \emph{Say it:} say \emph{nenapisikoḷḷu}
  \item \emph{Recall:} say the Kannada for to switch off, then the Kannada for to steal, then say \emph{nenapisikoḷḷu} again
\end{itemize}

\begin{tcolorbox}[breakable,colback=teal!4,colframe=teal!35!black,title={Before you move on}]
What does \kn{ನೆನಪಿಸಿಕೊಳ್ಳು} mean? (\textquotedblleft{}To recall, to remember\textquotedblright{}.) Say it once more.
\end{tcolorbox}
\section[\texorpdfstring{\kn{ಉರುಳು} (uruḷu)}{ಉರುಳು (uruḷu)}]{\kn{ಉರುಳು} (uruḷu) — to roll}

\label{lesson:KA-C165-urulu}



\begin{quote}
Before the new one: say the Kannada for to steal, then the Kannada for to recall.
\end{quote}

\subsection*{You'll want to know: \kn{ಉರುಳು}}
\textbf{\kn{ಉರುಳು}} — \emph{uruḷu} — \textquotedblleft{}to roll\textquotedblright{}.

\begin{grammarlens}[title={What you've built}]
One more word for this chapter.
\end{grammarlens}

\subsection*{Your turn}
\begin{itemize}
\raggedright
  \item \emph{Say it:} \emph{uruḷu}
  \item \emph{Say it:} \emph{uruḷu}, once more
  \item \emph{Say it:} say \emph{uruḷu}
  \item \emph{Recall:} say the Kannada for to steal, then the Kannada for to recall, then say \emph{uruḷu} again
\end{itemize}

\begin{tcolorbox}[breakable,colback=teal!4,colframe=teal!35!black,title={Before you move on}]
What does \kn{ಉರುಳು} mean? (\textquotedblleft{}To roll\textquotedblright{}.) Say it once more.
\end{tcolorbox}
\section[\texorpdfstring{\kn{ಉಳಿ} (uḷi)}{ಉಳಿ (uḷi)}]{\kn{ಉಳಿ} (uḷi) — to remain, to be left over}

\label{lesson:KA-C165-uli}



\begin{quote}
Before the new one: say the Kannada for to recall, then the Kannada for to roll.
\end{quote}

\subsection*{You'll want to know: \kn{ಉಳಿ}}
\textbf{\kn{ಉಳಿ}} — \emph{uḷi} — \textquotedblleft{}to remain, to be left over\textquotedblright{}.

\begin{grammarlens}[title={What you've built}]
One more word for this chapter.
\end{grammarlens}

\subsection*{Your turn}
\begin{itemize}
\raggedright
  \item \emph{Say it:} \emph{uḷi}
  \item \emph{Say it:} \emph{uḷi}, once more
  \item \emph{Say it:} say \emph{uḷi}
  \item \emph{Recall:} say the Kannada for to recall, then the Kannada for to roll, then say \emph{uḷi} again
\end{itemize}

\begin{tcolorbox}[breakable,colback=teal!4,colframe=teal!35!black,title={Before you move on}]
What does \kn{ಉಳಿ} mean? (\textquotedblleft{}To remain, to be left over\textquotedblright{}.) Say it once more.
\end{tcolorbox}
\section[\texorpdfstring{\kn{ಕಾಯಿಸು} (kāyisu)}{ಕಾಯಿಸು (kāyisu)}]{\kn{ಕಾಯಿಸು} (kāyisu) — to heat (milk, water)}

\label{lesson:KA-C165-kayisu}



\begin{quote}
Before the new one: say the Kannada for to roll, then the Kannada for to remain.
\end{quote}

\subsection*{You'll want to know: \kn{ಕಾಯಿಸು}}
\textbf{\kn{ಕಾಯಿಸು}} — \emph{kāyisu} — \textquotedblleft{}to heat (milk, water)\textquotedblright{}.

\begin{grammarlens}[title={What you've built}]
One more word for this chapter.
\end{grammarlens}

\subsection*{Your turn}
\begin{itemize}
\raggedright
  \item \emph{Say it:} \emph{kāyisu}
  \item \emph{Say it:} \emph{kāyisu}, once more
  \item \emph{Say it:} say \emph{kāyisu}
  \item \emph{Recall:} say the Kannada for to roll, then the Kannada for to remain, then say \emph{kāyisu} again
\end{itemize}

\begin{tcolorbox}[breakable,colback=teal!4,colframe=teal!35!black,title={Before you move on}]
What does \kn{ಕಾಯಿಸು} mean? (\textquotedblleft{}To heat (milk, water)\textquotedblright{}.) Say it once more.
\end{tcolorbox}
